\section{Reconstruction and further directions}\label{sec:further}

Theorems~\ref{thm:main} and~\ref{thm:integrated} distinguish the
information needed for integral reconstruction from that needed for a
leading spectral coefficient. The split-quintic construction determines
its global exterior and cohomological attachments. These recover the
displayed integral rings, specialization maps and vanishing lattice.
For the spectral law, the surviving polarization, the restriction of
the lost divisor class and its primitive intersection norm suffice
once the genuine geometric family satisfies the analytic hypotheses.
The common restriction map identifies which global classes fail to
extend over the smoothing and how they contribute to its small spectrum.
This is the specific connection between topology and geometry proved
here. It does not invert the spectrum to recover the whole smooth space.

The K3 reconstruction results
\cite{wuebben-reconstruction,wuebben-attachments} separate local models
from the global complement and its attachment information. The present
example makes that distinction outside the K3 bundle-moduli setting.
An extension to broader classes of threefolds would require a method for
calculating the global data from their algebraic degenerations, rather
than leaving an unknown complement ring among the conclusions. The
objects reconstructed here are the compact resolved, singular and
smoothed threefolds at the level of integral cohomology and the specified
maps; the cohomology ring of the complex four-dimensional total family
is a separate question.

Theorem~\ref{thm:matrix} determines the complete degree-two cluster for
any fixed number of lost directions. Its mass removes all surviving
harmonic classes, and its lower bound applies on every limiting spectral
subspace. Corollary~\ref{cor:eight-node} uses the compact construction
and exceptional data of \cite[Appendix~A, Proposition~A.1]{wuebben-chambers}
to obtain two distinct coefficients on that earlier family. This is a
direct geometric dependency within the programme. The relation to the
K3 attachment results remains methodological: their local models and
specialized reconstruction theorems are not premises of the conifold
spectral proof.

Corollary~\ref{cor:x9} supplies a direct connection to the separate
Calabi--Yau smoothing programme. CY~VI's original-side construction
provides $X'_9$, its projective small resolution and four exceptional
circuits. The nodal specialization of CY~V provides a full-support
deformation tangent and smooth deformation base. After extending the
specified ample class, the global intersections and the full surviving
kernel determine three explicit coefficients for the actual Ricci-flat
spectrum. This application uses earlier theorems as geometric premises.
The mixed-cone criterion and the mirror relation-lattice theorem have
different scopes; neither is silently reinterpreted as a nodal spectral
result. The curve-class map and its divisor-evaluation transpose explain
the common structure, while the quotient harmonic norm is additional
global information.

The finite-rank theorem gives leading limits for each fixed family.
It does not give relative errors uniform as positive generalized
eigenvalues approach zero in a varying family, nor does its full-kernel
projection provide the sharper one-sided remainder proved in rank one.

The present simple smoothing family has fixed nonzero leading node-scale
ratios. Allowing different vanishing orders or ratios approaching zero
requires actual-metric estimates uniform on that enlarged parameter domain.
The separate fixed-cap formula \eqref{eq:scaled-capacity} does not supply
those global estimates. On the split-quintic resolution, equality of all
exceptional curve classes also prevents independent exceptional areas.

The small two-form mode has an exact three-form partner under $d$, and
Hodge duality supplies related modes in degrees three and four. The
arguments above do not exhaust the middle-degree small spectrum. Likewise,
the iterated limits in the lower bound give a leading asymptotic rather
than a matching second-order remainder.

Finally, the tangent construction is intrinsic to the underlying
Ricci-flat geometry. A separately specified bundle requires its own
coupled deformation, stability and local connection analysis. Its
endomorphism Laplacian and HYM evolution are different operators and
problems from the real Hodge two-form theorem proved here. The local
Hermitian--Einstein and confined-flow theorems
\cite{wuebben-local-he,wuebben-hym-phase} retain a fixed smooth
background. Applying their conclusions uniformly through a conifold
transition would require new control of the background-dependent
inverses, bundle deformation and possible curvature concentration.
The no-atom tangent curvature estimate in
Proposition~\ref{prop:tangent-metrics} does not provide this analysis for
arbitrary bundles. Independently separating node scales, uniform control
through further degeneration of the spectral matrix and general
gauge-bundle limits are therefore distinct further steps in the
reconstruction programme.
